\chapter{Results and Evaluation}
\label{ch:results}

% ─────────────────────────────────────────────────────────────────────────────
\section{Overview of Evaluation Approach}
\label{sec:results_overview}
% ─────────────────────────────────────────────────────────────────────────────

Given the constraints documented in Section~\ref{subsec:meth_limitations}
--- specifically $n = 4$ external participants, no completed A/B control
condition, and IRT parameters calibrated from a small pilot sample --- the
evaluation adopts a \textbf{five-track multi-method validity} framework rather
than a single inferential statistical test. Each track provides a distinct
category of evidence about system validity, and convergent findings across
tracks constitute stronger evidence than any single track alone. This approach
follows the recommendation of \textcite{woolf2010} that ITS evaluation in
small-scale deployments should combine algorithmic, behavioural, and
psychometric evidence rather than relying exclusively on effect size estimates
that lack statistical power at small sample sizes.

The five tracks are:

\begin{description}
  \item[Track 1 — Algorithm validation.] Evidence that the adaptive engine
    behaves as specified: $\hat{\theta}$ trajectories move in the expected
    direction, Bloom escalation events fire at the correct thresholds, and
    variant deduplication prevents item reuse.
  \item[Track 2 — Engagement data.] Behavioural evidence of learner
    engagement: session completion rates, XP trajectories, questions
    attempted per session, and level progression.
  \item[Track 3 — Usability.] System Usability Scale (SUS) scores from
    $n = 4$ participants compared against the $\geq 68$ above-average
    benchmark \parencite{bangor2008}.
  \item[Track 4 — Misconception analysis.] Qualitative evidence of the
    system's diagnostic capability: distractor frequency tables per KC,
    identifying systematic knowledge gaps that a non-adaptive system would
    not surface.
  \item[Track 5 — System robustness.] Operational evidence: uptime,
    \texttt{variant\_pool\_exhausted} event frequency, and absence of
    session state corruption across the evaluation period.
\end{description}

% ─────────────────────────────────────────────────────────────────────────────
\section{Track 1 --- Algorithm Validation}
\label{sec:results_algorithm}
% ─────────────────────────────────────────────────────────────────────────────

\subsection{Theta Trajectory Analysis}
\label{subsec:results_theta}

For each participant session, the EAP ability estimate $\hat{\theta}$ was
recorded before and after each response. A valid adaptive session should
produce a $\hat{\theta}$ trajectory that: (a) moves upward when the learner
answers correctly, (b) moves downward or stabilises when the learner answers
incorrectly, and (c) converges toward a stable estimate as the number of
responses increases.

\begin{table}[H]
\centering
\caption{Session-level $\hat{\theta}$ statistics across $n = 4$ participants.
\textit{[Data to be inserted from Supabase Interaction table query.]}}
\label{tab:theta_trajectories}
\renewcommand{\arraystretch}{1.25}
\begin{tabular}{lccccc}
\toprule
\textbf{Participant} & \textbf{$\hat{\theta}$ start} & \textbf{$\hat{\theta}$ end}
  & \textbf{$\Delta\hat{\theta}$} & \textbf{Questions} & \textbf{Accuracy} \\
\midrule
P1 & {[TBA]} & {[TBA]} & {[TBA]} & {[TBA]} & {[TBA]}\% \\
P2 & {[TBA]} & {[TBA]} & {[TBA]} & {[TBA]} & {[TBA]}\% \\
P3 & {[TBA]} & {[TBA]} & {[TBA]} & {[TBA]} & {[TBA]}\% \\
P4 & {[TBA]} & {[TBA]} & {[TBA]} & {[TBA]} & {[TBA]}\% \\
\midrule
\textit{Mean} & \textit{{[TBA]}} & \textit{{[TBA]}} & \textit{{[TBA]}}
  & \textit{{[TBA]}} & \textit{{[TBA]}\%} \\
\bottomrule
\end{tabular}
\end{table}

A positive mean $\Delta\hat{\theta}$ across all participants would constitute
evidence that the EAP estimator is correctly updating in response to learner
performance, and that the session is producing measurable within-session
learning gains as operationalised by the IRT ability scale.

\subsection{Bloom Escalation Events}
\label{subsec:results_bloom}

The \texttt{bloom\_escalation} analytics event is emitted when the adaptive
engine advances a learner's Bloom gate for a KC. Valid escalation events
should satisfy: $p(\text{Learned}) \geq 0.40$ at L1$\to$L2 transition, or
$p(\text{Learned}) \geq 0.60$ at L2$\to$L3 transition.

\begin{table}[H]
\centering
\caption{\texttt{bloom\_escalation} events recorded across evaluation sessions.
\textit{[Data to be inserted from Supabase AnalyticsEvent table query.]}}
\label{tab:bloom_escalation}
\renewcommand{\arraystretch}{1.25}
\begin{tabular}{lcccc}
\toprule
\textbf{Session} & \textbf{KC} & \textbf{Transition} &
  \textbf{$p(\text{Learned})$ at trigger} & \textbf{Threshold met?} \\
\midrule
{[TBA]} & {[TBA]} & L1 $\to$ L2 & {[TBA]} & {[TBA]} \\
{[TBA]} & {[TBA]} & L2 $\to$ L3 & {[TBA]} & {[TBA]} \\
\bottomrule
\end{tabular}
\end{table}

All escalation events should show $p(\text{Learned})$ values at or above the
specified threshold at the moment of escalation, confirming that the gate
logic is correctly applied. Any event with $p(\text{Learned})$ below the
threshold would indicate a bug in the escalation condition check.

\subsection{Variant Deduplication Verification}
\label{subsec:results_dedup}

Cross-session deduplication was verified by inspecting the
\texttt{UserVariantSeen} table for each participant across multiple sessions.
\textit{[Verification result to be inserted: confirm zero repeat variants
served per participant across sessions.]} The \texttt{variant\_pool\_exhausted}
event was recorded \textit{{[TBA]}} times across all evaluation sessions,
indicating that the eligible pool was exhausted \textit{{[TBA]}} times and the
fallback selection logic was invoked. This confirms that the deduplication
mechanism is functioning and that variant pool size is a relevant operational
constraint for longer-running deployments.

% ─────────────────────────────────────────────────────────────────────────────
\section{Track 2 --- Engagement Data}
\label{sec:results_engagement}
% ─────────────────────────────────────────────────────────────────────────────

Engagement was assessed via the behavioural data recorded in the
\texttt{Session}, \texttt{Interaction}, and \texttt{XPEvent} tables.
Table~\ref{tab:engagement} summarises key engagement metrics across the
$n = 4$ participant sessions.

\begin{table}[H]
\centering
\caption{Participant engagement metrics.
\textit{[Data to be inserted from Supabase admin dashboard.]}}
\label{tab:engagement}
\renewcommand{\arraystretch}{1.25}
\begin{tabularx}{\textwidth}{lXXXX}
\toprule
\textbf{Participant} & \textbf{Session duration (min)} & \textbf{Questions attempted}
  & \textbf{KCs attempted} & \textbf{XP earned} \\
\midrule
P1 & {[TBA]} & {[TBA]} & {[TBA]} & {[TBA]} \\
P2 & {[TBA]} & {[TBA]} & {[TBA]} & {[TBA]} \\
P3 & {[TBA]} & {[TBA]} & {[TBA]} & {[TBA]} \\
P4 & {[TBA]} & {[TBA]} & {[TBA]} & {[TBA]} \\
\midrule
\textit{Mean} & \textit{{[TBA]}} & \textit{{[TBA]}} & \textit{{[TBA]}} & \textit{{[TBA]}} \\
\bottomrule
\end{tabularx}
\end{table}

Session completion rate (participants who reached the session summary screen
rather than abandoning mid-session) was \textit{{[TBA]}}\% ($n = $ \textit{{[TBA]}}
of 4). Hint usage rates per participant are reported in
Table~\ref{tab:hint_usage}, providing evidence on the scaffolding system's
uptake.

\begin{table}[H]
\centering
\caption{Hint usage by level across evaluation sessions.
\textit{[Data to be inserted from Supabase Interaction table query:}
\texttt{SELECT hintLevelMax, COUNT(*) FROM Interaction GROUP BY hintLevelMax}.
\textit{]}}
\label{tab:hint_usage}
\renewcommand{\arraystretch}{1.25}
\begin{tabular}{lcccc}
\toprule
\textbf{Hint level} & \textbf{Interactions} & \textbf{\% of total}
  & \textbf{Avg accuracy following hint} & \textbf{Avg credit factor} \\
\midrule
No hint (cf = 1.00) & {[TBA]} & {[TBA]}\% & {[TBA]}\% & 1.00 \\
L1 conceptual (cf = 0.90) & {[TBA]} & {[TBA]}\% & {[TBA]}\% & 0.90 \\
L2 procedural (cf = 0.70) & {[TBA]} & {[TBA]}\% & {[TBA]}\% & 0.70 \\
L3 worked example (cf = 0.50) & {[TBA]} & {[TBA]}\% & {[TBA]}\% & 0.50 \\
\bottomrule
\end{tabular}
\end{table}

% ─────────────────────────────────────────────────────────────────────────────
\section{Track 3 --- Usability (SUS)}
\label{sec:results_sus}
% ─────────────────────────────────────────────────────────────────────────────

The System Usability Scale \parencite{brooke1996} was administered to all
$n = 4$ participants at session end. The SUS comprises 10 Likert-scale items
(1--5), scored to produce a composite score on a 0--100 scale.
\textcite{bangor2008} established that scores $\geq 68$ indicate above-average
usability and scores $\geq 85$ indicate excellent usability for the product
category.

\begin{table}[H]
\centering
\caption{System Usability Scale scores ($n = 4$).
\textit{[Scores to be inserted from SUS response sheets.]}}
\label{tab:sus}
\renewcommand{\arraystretch}{1.25}
\begin{tabular}{lcc}
\toprule
\textbf{Participant} & \textbf{SUS score} & \textbf{Grade} \\
\midrule
P1 & {[TBA]} & {[TBA]} \\
P2 & {[TBA]} & {[TBA]} \\
P3 & {[TBA]} & {[TBA]} \\
P4 & {[TBA]} & {[TBA]} \\
\midrule
\textbf{Mean} & \textbf{{[TBA]}} & \textbf{{[TBA]}} \\
\textit{Benchmark ($\geq 68$)} & \textit{68} & \textit{Above average} \\
\bottomrule
\end{tabular}
\end{table}

A mean SUS score $\geq 68$ would indicate that the adaptive system does not
impose usability costs sufficient to impede learning, addressing the concern
raised by \textcite{bangor2008} that adaptive systems typically score 15--20
points below conventional interfaces. A score $\geq 85$ would indicate that
the onboarding tour, transparent BKT visualisation, and progressive hint
disclosure are effective usability interventions.

Qualitative SUS item-level analysis identifies which specific usability
dimensions contribute most to (or detract most from) the composite score.
Items 1, 3, 5, 7, and 9 are positive-direction (system confidence, ease of
learning, integration); items 2, 4, 6, 8, and 10 are negative-direction
(complexity, awkwardness, variability). \textit{[Item-level analysis to be
inserted.]}

% ─────────────────────────────────────────────────────────────────────────────
\section{Track 4 --- Misconception Analysis}
\label{sec:results_misconceptions}
% ─────────────────────────────────────────────────────────────────────────────

The \texttt{GET /api/analytics/error-patterns} route aggregates distractor
selection frequencies across all interactions for each KC. This enables
identification of systematic knowledge gaps beyond simple accuracy metrics:
a learner who consistently selects a specific wrong answer is exhibiting a
misconception, not random error.

\begin{table}[H]
\centering
\caption{Top distractor selections per KC across all evaluation sessions.
\textit{[Data to be inserted from error-patterns API response or direct
Supabase query against Interaction table.]}}
\label{tab:misconceptions}
\renewcommand{\arraystretch}{1.25}
\begin{tabularx}{\textwidth}{lXcc}
\toprule
\textbf{KC} & \textbf{Most common wrong answer} & \textbf{Frequency}
  & \textbf{\% of incorrect responses} \\
\midrule
{[TBA]} & {[TBA]} & {[TBA]} & {[TBA]}\% \\
{[TBA]} & {[TBA]} & {[TBA]} & {[TBA]}\% \\
{[TBA]} & {[TBA]} & {[TBA]} & {[TBA]}\% \\
\bottomrule
\end{tabularx}
\end{table}

The presence of high-frequency specific distractors (as opposed to uniformly
distributed errors) constitutes evidence that the distractor design is
functioning as intended: wrong answers are plausible misconceptions rather
than arbitrary foils, and the system's misconception tab surfaces genuine
diagnostic information for learner self-correction. \textit{[Qualitative
interpretation of identified misconception patterns to be inserted.]}

% ─────────────────────────────────────────────────────────────────────────────
\section{Track 5 --- System Robustness}
\label{sec:results_robustness}
% ─────────────────────────────────────────────────────────────────────────────

System robustness was assessed via operational metrics over the evaluation
period \textit{[dates: TBA]}.

\begin{table}[H]
\centering
\caption{System robustness metrics over evaluation period.
\textit{[Data to be inserted from Vercel deployment logs and Supabase
AnalyticsEvent table.]}}
\label{tab:robustness}
\renewcommand{\arraystretch}{1.25}
\begin{tabular}{lcc}
\toprule
\textbf{Metric} & \textbf{Value} & \textbf{Target} \\
\midrule
Uptime (Vercel deployment) & {[TBA]}\% & $\geq 99\%$ \\
Session state corruption events & {[TBA]} & 0 \\
\texttt{variant\_pool\_exhausted} events & {[TBA]} & Documented (fallback active) \\
Authentication failures (non-user-error) & {[TBA]} & 0 \\
RLS policy violation attempts blocked & {[TBA]} & All blocked \\
Average API response time (next-question) & {[TBA]} ms & $< 500$ ms \\
\bottomrule
\end{tabular}
\end{table}

Zero session state corruption events would confirm that the JSON serialisation
of \texttt{kcStates} and the atomic Prisma write in the answer route correctly
protect BKT state across concurrent requests. Any RLS policy violation attempts
blocked (reported in Supabase auth logs) would confirm that the data isolation
layer is functioning as a genuine security control, not merely as a
documentation claim.

% ─────────────────────────────────────────────────────────────────────────────
\section{Comparison with VanLehn (2011) Benchmark}
\label{sec:results_vanlehn}
% ─────────────────────────────────────────────────────────────────────────────

\textcite{vanlehn2011} establishes $d = 0.76$ as the expected effect size for
step-based ITS systems relative to no-tutoring control, across 60 controlled
studies. This project cannot compute a directly comparable effect size given
the absence of a no-tutoring control condition and the small sample ($n = 4$).
However, within-session $\Delta\hat{\theta}$ provides a proxy measure of
learning gain that can be interpreted against the IRT ability scale.

If participants' mean $\hat{\theta}$ increases by at least $0.50$ logits
across a session --- corresponding roughly to a one-half standard deviation
improvement in latent ability as estimated by the EAP --- this constitutes
evidence consistent with the magnitude of learning gains reported by VanLehn
for step-based systems, acknowledging that the comparison is approximate and
that $\hat{\theta}$ and effect size $d$ are not directly equivalent metrics.

\begin{table}[H]
\centering
\caption{Within-session learning gain compared against VanLehn benchmark.
\textit{[To be completed once session data is available.]}}
\label{tab:vanlehn_comparison}
\renewcommand{\arraystretch}{1.25}
\begin{tabular}{lcc}
\toprule
\textbf{Measure} & \textbf{This study} & \textbf{VanLehn (2011)} \\
\midrule
Comparison modality & Within-session $\Delta\hat{\theta}$ & Effect size $d$ \\
Sample & $n = 4$ (convenience) & 60 studies (controlled) \\
Mean learning gain & {[TBA]} logits & $d = 0.76$ \\
Inference possible? & No (descriptive only) & Yes (meta-analytic) \\
\bottomrule
\end{tabular}
\end{table}

The absence of inferential statistics is acknowledged as the primary
limitation of this evaluation. The comparison with VanLehn's benchmark is
offered as a contextualising reference, not as a statistical equivalence claim.
Chapter~\ref{ch:discussion} discusses this limitation in full and identifies
the study design required to produce a comparable effect size estimate.

% ─────────────────────────────────────────────────────────────────────────────
\section{Summary of Evaluation Findings}
\label{sec:results_summary}
% ─────────────────────────────────────────────────────────────────────────────

\begin{table}[H]
\centering
\caption{Five-track evaluation summary.}
\label{tab:eval_summary}
\renewcommand{\arraystretch}{1.25}
\begin{tabularx}{\textwidth}{llXl}
\toprule
\textbf{Track} & \textbf{Method} & \textbf{Finding} & \textbf{Status} \\
\midrule
1 & Algorithm validation &
  $\hat{\theta}$ trajectories, Bloom escalation events, dedup verification &
  \textit{{[TBA]}} \\
2 & Engagement data &
  Session completion, hint usage, XP trajectories &
  \textit{{[TBA]}} \\
3 & SUS usability &
  Mean SUS vs $\geq 68$ benchmark &
  \textit{{[TBA]}} \\
4 & Misconception analysis &
  Distractor frequency tables per KC &
  \textit{{[TBA]}} \\
5 & System robustness &
  Uptime, zero state corruption, RLS enforcement &
  \textit{{[TBA]}} \\
\bottomrule
\end{tabularx}
\end{table}

\noindent\textit{[This section will be completed with a convergent narrative
summary once all five tracks have been populated with empirical data from the
Supabase admin dashboard and SUS response sheets.]}
